% Image credits for the photographs and AI illustrations of Book 1.
% Machine-readable license records live in images/book1/CREDITS.md; keep the
% two files in sync. The Book 1 writer owns this file: add one \item per
% photograph (inside the itemize, by chapter) as photographs land, and remove
% the "(none yet)" item with the first one.
\chapter*{Image Credits}

The drawings, diagrams and charts in this book are original; charts drawn
from published data name their source in the caption or in this list.
Photographs are used under their respective free licenses, with thanks to
their authors; full source links and license URLs for every photograph are
listed in \texttt{images/book1/CREDITS.md} in the book's repository.

\bigskip
Alongside the photographs and the original drawings, some figures are
AI-generated illustrations, produced with OpenAI image generation from
prompts written by the book's editors, and reviewed for accuracy before
inclusion. They show generic scenes only, never a real person, a real
product or a real brand. The full prompt for every generated image is
recorded in \texttt{images/book1/ai/PROMPTS.md} in the repository.

\bigskip
\noindent Photographs, by chapter:
\begin{itemize}\small
  \item \textit{What Is a Computer?}: ENIAC programmed by Glen Beck and Betty
    Snyder, U.S. Army photograph, public domain (Wikimedia Commons).
  \item \textit{Typing with Ten Fingers}: a Sholes \& Glidden typewriter made
    by Remington, Wisconsin Historical Museum, photograph by Daderot, CC0
    (Wikimedia Commons).
  \item \textit{Inside the Computer}: part of a motherboard, photograph by
    Davidson F, CC0 (Wikimedia Commons).
  \item \textit{Memory and Storage}: an opened hard disk, photograph by
    Zzubnik, public domain; the disk mechanism of an IBM 350 at the Computer
    History Museum, photograph by Arnold Reinhold, CC BY-SA 3.0
    (both Wikimedia Commons).
  \item \textit{Networks and the Internet: A First Look}: the first web server,
    a NeXT computer at CERN, photograph by Coolcaesar, CC BY-SA 3.0
    (Wikimedia Commons).
  \item \textit{Repeating with \texttt{for}: Turtle Drawings}: a Valiant floor
    turtle robot and its drawing, image by Valiant Technology Ltd., CC BY-SA
    3.0 (Wikimedia Commons).
  \item \textit{How a Processor Runs a Program}: a processor in its socket,
    photograph by Capitalistcakes, CC0; knobs of an ENIAC function table,
    photograph by Logg Tandy, CC BY 4.0 (both Wikimedia Commons).
  \item \textit{Finding and Fixing Bugs}: the 1947 log-book page with the
    moth, U.S. Naval Historical Center photograph NH 96566-KN, public domain
    (Wikimedia Commons).
  \item \textit{How the Internet Works}: an opened optical fibre cable,
    photograph by Asurnipal, CC BY-SA 4.0 (Wikimedia Commons).
  \item \textit{Logic Gates and Boolean Algebra}: a 1997 replica of the first
    transistor, photograph by a U.S. federal employee, public domain
    (Wikimedia Commons).
  \item \textit{Secret Codes: A First Look at Cryptography}: a cipher wheel of
    the Swiss Army, photograph by Rama, CC BY-SA 2.0 FR (Wikimedia Commons).
\end{itemize}
